% Dados didaticos sem inventar orientacao, defesa ou catalogacao.
\titulo{A política dos números: construção, visibilidade e mobilização estatal de indicadores de violência letal no Brasil}
\subtitulo{}
\titulotraduzido{The politics of numbers: construction, visibility and state mobilization of lethal violence indicators in Brazil}
\autor{Rafael Carvalho dos Santos}
\autorficha{SANTOS, Rafael Carvalho dos}
\orientador{Orientação não informada neste exercício}
\coorientador{}
\curso{Administração Pública}
\programa{Administração Pública}
\area{}
\linhapesquisa{}
\cidade{Brasília}
\local{Brasília-DF}
\data{2026}
\datadefesacompleta{Exercício didático, sem defesa ou aprovação acadêmica}
\cutter{Não se aplica ao exercício}
\numeropaginas{Não se aplica à catalogação deste exercício}
\cdd{Não se aplica ao exercício}
\assuntos{Indicadores. Violência letal. Comunicação institucional. Mobilização estatal.}
\palavraschave{Indicadores. Violência letal. Visibilidade midiática. Mobilização estatal.}
\keywords{Indicators. Lethal violence. Media visibility. State mobilization.}
\preambulo{Esqueleto didático de dissertação elaborado para o exercício ia-6.1 da disciplina Agentes de IA, com base no projeto do Mestrado Profissional em Administração Pública do IDP. O documento testa o template LaTeX e não apresenta resultados empíricos nem constitui dissertação concluída.}
